\section{What remains for a transposition algorithm or impossibility proof}
\label{sec:roadmap}

The main conclusion is a uniform deficit for specified full return
geometries. It is stronger than a statement that one proposed code has
failed: every deterministic local computation on any geometry in the
stated class is covered, with no linearity assumption. Its scope is also
explicit. A bounded full-span partition, a Q2 core, and a bounded number
of exceptional families are additional geometric hypotheses. They have
not been derived for an arbitrary linear-time transposer.

\subsection{The nearest unresolved geometric case}

For one exceptional family, \cref{cor:one-exception} completes the
deficit-versus-payment argument. With two exceptional families, a
limiting inverse-square hop profile is compatible with
\cref{thm:necessary-hop-tail}. The next concrete problem is to determine
whether the remaining complete continuation can satisfy all the
necessary conditions simultaneously.

There are three precise ways to exclude a common native class in this
case. First, bound the phase-and-column mean fractional cover by a
common function $f(R)=o(R^2)$, uniformly in dimension and every finite
first-$J$ selection with $J\ge R$. Second, bound the mean number of
globally maximal newer seeds, at each fixed protected scale $R$ in a
common unbounded set, by a function $g_R(J)$ with
$\liminf_J g_R(J)/(J+1)^2=0$. Third, give the raw metric covers with
the three joint limits in \cref{eq:metric-sharing-limits}.
These are respectively the criteria in
\cref{cor:uniform-cover,thm:uniform-sharing,prop:metric-sharing}.
They are alternatives; none is currently supplied for every admissible
two-exception geometry with an inverse-square envelope.

A failed attempt should be checked against the exact hypothesis it
tests. \Cref{prop:serial-quadratic-cover} defeats a universal pointwise
linear cover estimate based only on serial exceptional spans and
exceptional-only variation tails. It does not defeat a phase average,
a common inverse-square envelope, or shared raw revelation across all
outputs. It also belongs to a class already excluded by the finite
skeleton corollary. Using it as evidence against all three uniform
criteria would therefore draw a conclusion the example does not support.

Conversely, a candidate near-full sequence in the inverse-square class
must meet the necessary reciprocal-cover, complete-observable movement,
and global-continuation-density bounds of
\cref{thm:necessary-cover,cor:necessary-exposure,cor:sharing-density},
with all limits taken along the same joint profile. Producing such a
geometry would test the sharpness of these certificates. It would not
by itself supply a correct computation on the geometry.
Even without a size envelope, \cref{cor:sharing-density} forces
$g^{\rm low}_{J,R}/(J+1)\to\infty$ at every sufficiently large fixed
$R$. A dimension-uniform linear seed bound would therefore already
exclude a near-full sequence in that class.

\subsection{What a constructive advance must compute}

A recovery code becomes a transposition algorithm only after its
encoding, updates, reads and output order have been implemented within
the native time bound. The complete-source certificate charges one
such implementation; it does not construct one from an abstract code.
The coded-caching calibration makes this distinction concrete. Its
algebraic decoder exists, but its specified raw-query interface cannot
be represented by a fixed number of trajectories with uniformly bounded
normalized movement at all small fixed radius ratios. Changing that
interface or helper may evade the argument and must be analysed anew.

The causal-group result identifies another restricted shortcut that
cannot work: after a cheap helper is fixed, a locally autonomous
computation for fixed groups of $o(\log d)$ adjacent columns cannot
retain nearly all the input information. A successful algorithm can
avoid that premise by coupling groups during computation. The theorem
does not price arbitrary such coupling. The constructive target is an
explicit coupled update rule whose total native execution cost can be
bounded, including the cost of changing where its encoded information
is available. Final-output separation alone is not the forbidden
causal typing.

\subsection{What would decide (T)}

An affirmative answer requires one fixed finite-tape, finite-alphabet
machine, constants $c,d_0$, and a proof that it correctly transposes
every input of every dimension $d\ge d_0$ in at most $cd^2b(d)$ steps.
All preprocessing and recovery must be included. A family of
dimension-dependent machines, a source relation, or a storage code
without such an implementation does not meet this specification.

A negative answer must exclude every such fixed machine and constant.
One possible route is a reduction from an arbitrary candidate to
correctness sets and paid certificates with a strictly positive net
deficit. Such a reduction must account for the cost of selecting a
geometry in the adaptive case, and it must address geometries outside
the span and core hypotheses used here. Alternatively, a different
certificate may handle those geometries directly. The depth reductions
of the fourth and eighth papers and the complete-recovery restrictions
of the fifth provide complementary tools
\cite{FourthPaper2026,EighthPaper2026,DirectedRecovery2026}; this paper
does not prove that their hypotheses cover every remaining machine.

The formalization verifies the finite certificates and their stated
uniform consequences. It does not promote a geometric criterion into
a universal machine property. The following appendices record that
boundary, the dependencies and the verification commands.
